\subsection{Conditional Trajectory Diffusion for Dynamic Constraint Satisfaction}
\label{sec:discussion:diffusion}

\cref{ch:generation,ch:payload} shifted from explicit kinodynamic search to learned generative representations of dynamic feasibility.
Rather than constructing a trajectory online through repeated propagation, collision checking, or constraint evaluation, the approach reframes motion generation as conditional sampling: a model learns the distribution of trajectories that satisfy dynamic constraints under specified task conditions, then samples from that distribution at inference time.

The primary instantiation of this idea is payload-conditioned trajectory diffusion.
As payload mass increases, the feasible motion set changes because joint torque, acceleration, and inertial constraints become more restrictive.
Instead of enforcing these constraints through runtime checking, the diffusion model conditions directly on payload mass and learns to generate joint-space trajectories whose positions, velocities, and accelerations are consistent with the corresponding dynamic limits.
The result is near-constant-time trajectory generation for super-nominal payload manipulation, maintaining substantial workspace accessibility even at three times the robot's rated payload capacity.

The case study on failure-induced embodiments demonstrates that the same conditioning principle extends beyond payload variation.
In \textit{DEFT}, actuation failures are represented as changes in embodiment, and the diffusion model conditions on the resulting failure configuration when generating trajectories.
This allows a single model to adapt across multiple joint-failure conditions and motion primitives, treating failure not as an exception requiring a separate controller, but as another structured input to the generative process.

Together, the payload and failure results show that dynamic constraints can be encoded implicitly through conditioning when the conditioning signal captures the physical factors that shape feasibility.
This marks a conceptual transition in the thesis: dynamic feasibility is no longer evaluated only by a planner or simulator, but represented as a learned conditional distribution over trajectories.

% \subsection{Conditional Trajectory Diffusion for Dynamic Constraint Satisfaction}
% \label{sec:discussion:diffusion}

% The second thematic area of the thesis (\cref{ch:generation,ch:payload}) reframes motion planning as conditional sampling: rather than searching or optimizing for a trajectory that satisfies constraints, a generative model learns to sample from the distribution of constraint-satisfying trajectories conditioned on task-level properties.

% For super-nominal payload manipulation, the conditioning signal is payload mass.
% The model is trained on trajectories generated across the full range of payload-dependent dynamic constraints using multiple planners operating in varied environments, so the diffusion process embeds constraint satisfaction in generation rather than enforcing it post-hoc.
% At inference time, the model generates joint-space trajectories---position, velocity, and acceleration---in approximately 10\,ms, maintaining 67.6\% workspace accessibility at three times the robot's nominal payload rating without runtime constraint checking.
% This establishes that dynamic constraints can be implicitly encoded in a generative model through conditioning alone.

% \textit{DEFT} extends this conditioning principle to joint failures.
% Rather than treating each failure configuration as a special case requiring replanning or policy switching, \textit{DEFT} conditions the same diffusion architecture on an embodiment representation encoding the failure configuration.
% A single model generalizes across multiple failure types and motion primitives.
% Real-world experiments demonstrate high task success under previously unseen multi-joint failure combinations with strict joint-limit compliance.
% Together, the payload and failure results show that any structured property of the task or robot state can serve as a conditioning signal for dynamics-constrained trajectory generation.

\subsection{Task-Agnostic Seeding for Constrained Trajectory Optimization}
\label{sec:discussion:seeding}

\cref{ch:seeding} extended the generative view of dynamic feasibility to constrained trajectory optimization.
While trajectory optimization provides a powerful framework for enforcing complex kinematic and dynamic constraints, its performance depends heavily on initialization.
Poor seeds can cause solvers to converge slowly, become trapped in infeasible regions, or fail entirely, especially when the feasible set is narrow or shaped by higher-order dynamic constraints.

The proposed approach uses diffusion not as a final trajectory generator, but as an initialization mechanism.
Instead of conditioning on a task-specific scalar such as payload mass or an embodiment descriptor such as a failure configuration, the model conditions on informative samples of robot state---joint positions, velocities, accelerations, and torques---that satisfy the relevant constraints.
These samples provide a task-agnostic representation of the constraint manifold, allowing the model to generate trajectory seeds that begin near dynamically feasible regions of the search space.

% The central contribution is therefore not a new optimizer, but a general interface between learned generation and constrained optimization.
By translating arbitrary dynamic constraints into state-sample conditioning, this general interface between learned generation and constrained optimization reduces the amount of task-specific engineering required to initialize optimization-based planners.
Across diverse manipulation tasks, this produces seeds that improve solver convergence and feasibility while preserving the ability of the downstream optimizer to enforce constraints explicitly.

In the broader arc of the thesis, task-agnostic seeding represents the most abstract representation of dynamic feasibility developed here.
Where earlier chapters evaluate dynamics online, reuse precomputed motion structure, or learn task-conditioned trajectory distributions, this approach learns to generate useful initializations from samples of the feasible set itself.
It suggests a general strategy for combining learned priors with optimization: use learning to move close to the constraint manifold, then use optimization to refine and certify the final motion.

% \subsection{Task-Agnostic Seeding for Constrained Trajectory Optimization}
% \label{sec:discussion:seeding}

% The final contribution (\cref{ch:seeding}) addresses the initialization sensitivity of constrained trajectory optimization.
% Existing seeding methods are either task-specific or limited to geometric constraints; when dynamics constraints enter the optimization, poor initialization causes convergence failure or infeasible solutions.
% The proposed approach conditions a diffusion model on samples of robot state---joint angles, velocities, accelerations, and torques---that satisfy task-level dynamic constraints, producing near-feasible trajectory seeds without task-specific engineering.
% By initializing the optimizer near the constraint manifold, the approach substantially improves convergence and feasibility across diverse tasks.
% The deeper contribution is architectural: a single conditioning mechanism, agnostic to the specific task or constraint type, can serve as the initialization front-end for any dynamics-constrained trajectory optimizer.